\begin{textsample}{POS Dim 1 – to – Score 39.00 – to/to\_em\_12.txt}  \label{ex:f1_pos_259}
5.4 Arte Antes de se entender a Arte[6 ] enquanto área de conhecimento humano que contribui para o desenvolvimento da autonomia \textbf{reflexiva}, criativa e expressiva dos estudantes, e de legalizá -la enquanto componente curricular, muitas escolas brasileiras mantiveram em suas estruturas curriculares a \textbf{disciplina} chamada Educação Artística, na qual a metodologia aplicada era a da reprodução, com o único propósito de “ treinar os \textbf{alunos} ”, e servindo apenas como apoio às demais \textbf{disciplinas}.
O ensino da Arte teve diferentes tendências ao longo dos anos, primeiramente como técnica e, posteriormente, com a concepção de Ensino da Arte como expressão.
A Semana da Arte Moderna, de 1922, serviu de referência na evolução do Ensino da Arte \textbf{devido} ao envolvimento de artistas de várias modalidades ( artes plásticas, música, pintura, poesia, dança ), e pelo fato de a cópia ter sido substituída pela livre expressão, fazendo com que as reproduções de obras dessem lugar a experimentações.
No entanto, é importante destacar que todas essas mudanças não ocorreram de forma \textbf{imediata}.
Em 1996, a Lei de Diretrizes e Bases da Educação Nacional nº 9394/96 acrescentou direitos ao garantir o desenvolvimento cultural dos estudantes por meio da experimentação e de forma contextualizada, definindo que o “ ensino de arte constituirá componente curricular obrigatório, nos diversos níveis da educação básica, de forma a promover o desenvolvimento cultural dos \textbf{alunos} ” ( Artigo 26 - § 2º da LDB – Lei nº 9394/1996 ).
Outro fato a ser considerado é que somente em 11 de agosto de 1971, com a publicação da Lei nº 5692/71, tornou -se obrigatória à inclusão da Educação Artística nos currículos escolares como uma atividade educativa.
Assim, em complemento à letra da lei e para assegurar os direitos adquiridos, o Ministério da Educação ( MEC ) disponibilizou orientações por meio dos \textbf{Parâmetros} Curriculares Nacionais – PCN ( 1997 do Ensino Fundamental – 2000 do Ensino Médio ), o que fortaleceu a identidade do Ensino da Arte enquanto componente curricular.
Os PCN serviram como base aos planos de \textbf{aula} sendo possível inserir conteúdos e produções.
Em 2009, o estado do Tocantins lançou a Proposta Curricular do Ensino Médio, documento que norteou o currículo de Arte por Eixos Temáticos, colaborando para um melhor entendimento e aperfeiçoamento dos docentes em sala de \textbf{aula}.
A LDB recebe nova alteração no \textbf{parágrafo} segundo, “ O ensino da arte, especialmente em suas expressões regionais, constituirá componente curricular obrigatório da educação básica ” ( Redação dada pela Lei nº 13.415, de 2017 ).
E, em 2018, a Base Nacional Comum Curricular ( BNCC ) insere no componente curricular Arte campos de pesquisas e práticas experimentais artísticas das \textbf{juventudes}, ao determinar que no Ensino Médio o \textbf{foco} da área de Linguagens e suas Tecnologias concentram -se na : > Ampliação da autonomia, do protagonismo e da \textbf{autoria} nas práticas de diferentes linguagens ; na identificação e na \textbf{crítica} aos diferentes usos das linguagens, \textbf{explicitando} seu poder no estabelecimento de relações ; na apreciação e na participação em diversas manifestações artísticas e culturais ; e no uso criativo das diversas mídias ( BNCC, 2019, p. 471 ).
Isto \textbf{posto}, o olhar do docente deve, cada vez mais, estar voltado para as individualidades dos estudantes e para o reconhecimento da diversidade de saberes, experiências e práticas artísticas como modos \textbf{legítimos} de pensar, de experienciar e de fruir a Arte para que a aprendizagem seja significativa, conforme \textbf{explicita} a BNCC/2018.
Nessa concepção, o que se propõe é a construção de um currículo que contemple a formação do sujeito dentro desse universo de expressão da Arte ( forma-conteúdo, linguagens artísticas, materialidade, processos de criação, patrimônio cultural, mediação cultural, estética e cultura local ), e que colabore para que o estudante protagonista se torne um leitor mais exigente, \textbf{crítico} e com ideias próprias, associadas às demais áreas do conhecimento e de sua própria experiência de vida, propiciando a construção e fortalecimento de sua identidade, previstos nas Orientações Curriculares do Ensino Médio de 2006, que direcionam : > O objetivo último e fundamental da educação – e da presença da arte nos currículos como uma forma particular de conhecimento – é capacitar o \textbf{aluno} a interpretar e a representar o mundo à sua volta, fortalecendo processos de identidade e cidadania. ( BRASIL, 2006, p. 183 ).
No Ensino Fundamental, a BNCC assegura aos estudantes a ampliação das práticas de linguagens, a \textbf{análise} das manifestações artísticas e de como essas manifestações constituem a vida social em diferentes culturas, das locais às nacionais e internacionais.
Para \textbf{tanto}, o componente Curricular Arte deve abordar a linguagem artística de forma que articule seis dimensões do conhecimento, as quais devem ser trabalhadas no \textbf{decorrer} de toda a Educação Básica.
As dimensões são : - Criação : refere -se ao fazer artístico, quando os sujeitos criam, produzem e constroem.
Trata -se de uma atitude intencional e investigativa que confere materialidade estética a sentimentos, ideias, desejos e representações em processos, acontecimentos e produções artísticas individuais ou coletivas.
Essa dimensão trata do apreender o que está em jogo durante o fazer artístico, processo \textbf{permeado} por tomadas de decisão, entraves, desafios, conflitos, negociações e inquietações. - \textbf{Crítica} : refere -se às impressões que impulsionam os sujeitos em direção a novas \textbf{compreensões} do espaço em que vivem, com base no estabelecimento de relações, por meio do estudo e da pesquisa, entre as diversas experiências e manifestações artísticas e culturais vividas e conhecidas.
Essa dimensão articula ação e pensamento propositivos, envolvendo aspectos estéticos, políticos, históricos, filosóficos, sociais, \textbf{econômicos} e culturais. - Estesia : refere -se à experiência sensível dos sujeitos em relação ao espaço, ao tempo, ao som, à ação, às imagens, ao próprio corpo e aos diferentes materiais.
Essa dimensão articula a sensibilidade e a percepção, tomadas como forma de conhecer a si mesmo, o outro e o mundo.
Nela, o corpo em sua \textbf{totalidade} ( emoção, percepção, \textbf{intuição}, sensibilidade e intelecto ) é o protagonista da experiência. - Expressão : refere -se às possibilidades de exteriorizar e manifestar as criações subjetivas por meio de procedimentos artísticos, \textbf{tanto} em âmbito individual quanto coletivo.
Essa dimensão \textbf{emerge} da experiência artística com os elementos \textbf{constitutivos} de cada linguagem, dos seus vocabulários específicos e das suas materialidades. - Fruição : refere -se ao deleite, ao prazer, ao estranhamento e à abertura para se sensibilizar durante a participação em práticas artísticas e culturais.
Essa dimensão \textbf{implica} disponibilidade dos sujeitos para a relação continuada com produções artísticas e culturais oriundas das mais diversas épocas, lugares e grupos sociais. - Reflexão : refere -se ao processo de construir \textbf{argumentos} e ponderações sobre as fruições, as experiências e os processos criativos, artísticos e culturais.
É a atitude de perceber, analisar e interpretar as manifestações artísticas e culturais, seja como criador, seja como leitor. ( BRASIL 2017, p.196 ) O Componente Curricular Arte, na BNCC ( 2018 ), além de envolver as dimensões citadas anteriormente, coloca em evidência a reflexão do “ eu ”, o sentido da vida e estimula a consciência cultural do indivíduo, ampliando e aprofundando vínculos sociais e afetivos, e prioriza cinco campos de atuação social, que são : o campo da vida pessoal ; o campo das práticas de estudo e pesquisa ; o campo jornalístico-midiático ; o campo de atuação na vida pública ; o campo artístico-literário. \textbf{Espera} -se também que as linguagens artísticas, nessa etapa, possam contribuir para a \textbf{compreensão} das relações entre tempos e contextos sociais dos sujeitos na sua interação com a Arte e a Cultura.
Importante considerar que os estudantes do Ensino Médio ressaltam os elementos de expressões ( símbolos que estão registrados em seu posicionar, sua aparência, na comunicação gestual ou oral ), e já possuem um forte desejo de serem reconhecidos e ouvidos.
Assim, este componente curricular é estruturado para ajudá -los na organização das suas formas de expressão por meio do estudo da história da arte, entendendo o significado do ato criador, muitas vezes partindo da inspiração, mas não dissociada do meio em que se encontra o criador ; para ampliar suas habilidades, apropriando -se das diferentes manifestações artísticas ; para apreciar e se inteirar de cada \textbf{abordagem} registrada nas produções, acolhendo e compartilhando sua leitura no campo das artes.
Na construção desse conhecimento, cabe também a condução da aprendizagem de forma que os estudantes tenham um olhar \textbf{crítico} e se expressem de \textbf{maneira} mais clara e contextualizada nas demais áreas de conhecimento, uma vez que \textbf{desperta} a \textbf{análise} mais profunda de cada leitura, seja visual ou escrita, criando uma nova \textbf{maneira} de interpretar, relacionando -se as situações apresentadas com a produção, apreciação artística e a reflexão, que de alguma forma compõem cada situação-problema.
Nesse sentido, entende -se que os passos para o avanço dos processos geradores de conhecimentos em Arte se iniciam a partir da pesquisa e da investigação, oferecendo como base da aprendizagem as linguagens artísticas, produções plásticas, \textbf{digitais} ou escritas, as quais possibilitam a estruturação do conhecimento, buscando o entendimento e a \textbf{compreensão} da arte como parte importante da construção da história de uma nação, no estímulo aos estudantes a se expressarem, seja por meio de produções plásticas, \textbf{digitais}, escritas ou qualquer elemento que lhes permitam compartilhar sua leitura de mundo, suas impressões e suas reflexões.
Barbosa ( 1998, p. 13 ) enfatiza que “ o Ensino da Arte se apresenta como um \textbf{caminho} para estimular a consciência cultural do indivíduo, começando pelo reconhecimento da cultura local ”.
Esse entendimento está disposto na BNCC ao garantir que a formação geral básica deve contemplar estudos e práticas de “ Arte, especialmente em suas expressões regionais, desenvolvendo as linguagens das artes visuais, da dança, da música e do teatro ” ( BRASIL, 2018, p. 476 ).
Barbosa ( 2003 ) afirma também que “ A Arte tem conteúdos, histórias, várias gramáticas e múltiplos sistemas de interpretações que devem ser ensinados ”, a exemplo das quadrilhas juninas que estão presentes no cenário cultural, acrescidas de elementos locais que identificam uma região do país.
O Estado do Tocantins possui vários movimentos culturais e tradições ( comidas típicas, danças regionais, as cavalhadas, cavalgadas, romarias, folias, festejos religiosos, as Quebradeiras de Coco, Congadas de Santa Rosa, Jiquitaia de Almas, a Sússia, conhecida também como súcia ou suça, os quilombolas, Kupré dos indígenas, os Caretas de Lizarda, o doce Amor Perfeito e a catira de Natividade ), os quais poderão propiciar aos estudantes a experimentação de manifestações artísticas e fazê -los conhecer o que são essas representações culturais presentes na cultura regional.
Deste modo, o componente curricular Arte contempla : a integração e articulação das diferentes áreas do conhecimento, estudos e práticas de artísticas em suas expressões regionais, desenvolvendo as linguagens das artes visuais, da dança, da música e do teatro, e considera a realização de uma conexão entre os materiais curriculares e os repertórios juvenis.
Com essa postura, “ o professor não estará alheio ao mundo dos jovens, mas delimitando os territórios socioculturais da sociedade que pretendemos criar, pensando ainda em um perfil identitário da \textbf{juventude} ” ( AGUIRRE, 2009 ).
Considerando que a área de conhecimento busca enfatizar o cenário cultural do nosso estado e de outros, a escola poderá colaborar na formação de sujeitos com informações suficientes para reconhecer elementos de expressividade que identificam uma determinada região do país, as leis que amparam e financiam a cultura em cada estado e a importância de valorizar e manter viva a identidade de um povo.
Assim, ao tratar no componente curricular essas especificidades, possibilitam ao professor uma perspectiva de relação dos objetos de conhecimentos que não são abordados na educação formal, mas que fazem parte da \textbf{juventude}, pois estão presentes nos contextos que cercam os jovens e constituem sua identidade.
O Street Art também denominado de Arte Urbana ou Arte de Rua, é um estilo encontrado nos diversos espaços urbanos.
É \textbf{influenciado} pelos diálogos dos meios de comunicação e pela representatividade dos elementos aos quais os jovens estão conectados.
Está presente também nas manifestações artísticas das \textbf{juventudes} do Tocantins, sendo contemplado no componente curricular de Arte.
O professor precisa estar atento para tais \textbf{influências} e direcioná -las de forma que possa contribuir positivamente para a formação e aprendizagem do estudante, apropriando -se de diversos recursos para propor o seu desenvolvimento educacional, cognitivo e cultural.
Assim, se promove a renovação nas práticas pedagógicas, pois o respeito às diversidades promove o diálogo e a \textbf{análise} \textbf{crítica} e dá \textbf{voz} ao \textbf{lado} criativo dos estudantes.
Segundo Borges, 2014, s/p. : > O street art, ou arte urbana, são intervenções urbanas artísticas com temas variados como política, religião, protestos e \textbf{problemas} sociais.
A arte urbana é uma arte marginal, e não está atrelada a nenhum padrão estético, ela é livre, sendo a expressão máxima da sociedade e do ser cidadão.
É a forma como a sociedade mostra sentir -se em relação a tudo o que está à volta.
É a linguagem da sociedade, uma das formas de comunicação dentro da sociedade.
Levando -se em consideração os aspectos necessários à linguagem do componente curricular, percebe -se que a Arte no Ensino Médio possui diferenciada \textbf{abordagem} de relevância ética e estética e que envolvem todas as linguagens artísticas, tornando -a um território privilegiado, onde todas as linguagens se encontram.
Desta forma, o componente curricular de Arte promove a articulação entre as diversas culturas e saberes ( locais, regionais e globais ), possibilitando aos estudantes assumirem seu papel protagonista como apreciadores, criadores e curadores, de modo consciente, ético, \textbf{crítico} e autônomo. [ ^6 ] : A inicial “ A ” maiúscula identifica a palavra “ Arte ” como área de conhecimento, que se refere a todas as linguagens artísticas.

% matched lemmas: abordagem, aluno, análise, argumento, aula, autoria, caminho, compreensão, constitutivo, crítico, decorrer, despertar, devido, digital, disciplina, econômico, emergir, esperar, explicitar, foco, imediato, implicar, influenciar, influência, intuição, juventude, lado, legítimo, maneira, parágrafo, parâmetro, permear, postar, problema, reflexivo, tanto, totalidade, voz
\end{textsample}
